\section{Operational Deep Dives}

\subsection{Red Team Report Workflow}

After every capability upgrade, the red team produces a narrative report. It includes timeline, trigger evidence, and mitigation steps. The narrative is structured as (1) pre-incident context, (2) capability vector, (3) observed behavior, (4) hypothesized failure mode, (5) recommended controls.

The report is reviewed by the governance council and ingested into the Capability Report. When multiple reports cite the same failure mode, an `Issue Cluster` is created and assigned to a dedicated `Capability Owner`.

\subsection{Compliance and Transparency Reporting}

Anthropic publishes regular transparency briefs describing: model changes, test results, and outstanding concerns. Each brief is validated against the Long-term Benefit Trust board schedule.

The transparency workflow includes three sign-offs: research lead, safety lead, trustee. Any deviation before sign-off triggers a pause and a formal explanation to the board.

\subsection{Data Logging Strategy}

All instrumentation data (capability metrics, tool calls, overrides) are streamed into a multi-tenant warehouse with `[tag, timestamp, event\_id, payload]` records. The ingestion pipeline tags events with `project`, `agent\_version`, `ASL\_level`.

Raw logs are retained for 90 days; aggregated dashboards (alignment drift, tool success) are kept for 365 days. This allows retrospective audits long after a deployment.

\subsection{本章小结}

These deep dives illustrate how red teaming, transparency, and logging are not ad-hoc but woven into regular governance rituals. Each ritual produces artifacts that feed benchmarking, monitoring, and deployment guardrails.
